\chapter{PENGUJIAN DAN ANALISIS}
\label{sec:chap4_pengujian}

\section{Skenario Pengujian}

Pengujian dilakukan untuk memastikan seluruh komponen sistem dapat beroperasi sesuai dengan kebutuhan yang telah ditentukan pada tahap perancangan. Pengujian mencakup aspek komunikasi data, penyimpanan data, monitoring real-time, motion detection, notifikasi mobile, monitoring telemetri, serta performa server.

\section{Pengujian Koneksi Perangkat ESP32-CAM}

Pengujian ini dilakukan untuk memastikan perangkat ESP32-CAM dapat terhubung ke jaringan Wi-Fi dan melakukan komunikasi dengan broker MQTT.

\subsection{Prosedur Pengujian}

Langkah-langkah pengujian yang dilakukan adalah sebagai berikut:

\begin{enumerate}
	\item Menyalakan perangkat ESP32-CAM.
	\item Menghubungkan perangkat ke jaringan Wi-Fi.
	\item Melakukan koneksi ke broker MQTT.
	\item Mengamati status koneksi melalui dashboard monitoring.
\end{enumerate}

\subsection{Hasil Pengujian}

\begin{table}[H]
	\centering
	\caption{Hasil Pengujian Koneksi ESP32-CAM}
	\begin{tabular}{|c|c|c|}
		\hline
		No & Parameter & Hasil \\
		\hline
		1 & Koneksi Wi-Fi & Berhasil \\
		2 & Koneksi MQTT & Berhasil \\
		3 & Pengiriman Data & Berhasil \\
		4 & Reconnect Otomatis & Berhasil \\
		\hline
	\end{tabular}
\end{table}

Berdasarkan hasil pengujian, perangkat ESP32-CAM berhasil terhubung ke jaringan dan melakukan komunikasi dengan server melalui broker MQTT secara stabil.

\section{Pengujian Komunikasi MQTT}

Pengujian dilakukan untuk memastikan data yang dikirimkan oleh ESP32-CAM dapat diterima oleh server tanpa kehilangan data.

\subsection{Prosedur Pengujian}

\begin{enumerate}
	\item ESP32-CAM mengirimkan citra melalui MQTT.
	\item Broker EMQX menerima data.
	\item Laravel memproses data yang diterima.
	\item Data disimpan ke sistem penyimpanan.
\end{enumerate}

\subsection{Hasil Pengujian}

\begin{table}[H]
	\centering
	\caption{Hasil Pengujian MQTT}
	\begin{tabular}{|c|c|}
		\hline
		Parameter & Hasil \\
		\hline
		Publish Data & Berhasil \\
		Subscribe Data & Berhasil \\
		Retain Message & Berhasil \\
		QoS & Berhasil \\
		\hline
	\end{tabular}
\end{table}

Hasil pengujian menunjukkan bahwa komunikasi data menggunakan MQTT berjalan dengan baik dan mampu mendukung proses pengiriman data dari perangkat ke server.

\section{Pengujian Penyimpanan Data}

Pengujian dilakukan untuk memastikan data citra dan metadata dapat tersimpan dengan benar.

\subsection{Hasil Pengujian}

\begin{table}[H]
	\centering
	\caption{Hasil Pengujian Penyimpanan Data}
	\begin{tabular}{|c|c|}
		\hline
		Jenis Data & Hasil \\
		\hline
		Metadata Kamera & Berhasil \\
		Data Telemetri & Berhasil \\
		Notifikasi & Berhasil \\
		Citra Kamera & Berhasil \\
		\hline
	\end{tabular}
\end{table}

Data metadata berhasil tersimpan pada MySQL sedangkan file citra berhasil tersimpan pada MinIO Object Storage.

\section{Pengujian Website Real-Time}

Pengujian dilakukan untuk memastikan website mampu menerima pembaruan data secara real-time menggunakan WebSocket.

\subsection{Prosedur Pengujian}

\begin{enumerate}
	\item Mengaktifkan perangkat kamera.
	\item Mengirimkan data ke server.
	\item Mengamati perubahan data pada dashboard.
\end{enumerate}

\subsection{Hasil Pengujian}

\begin{table}[H]
	\centering
	\caption{Hasil Pengujian Website Real-Time}
	\begin{tabular}{|c|c|}
		\hline
		Fitur & Hasil \\
		\hline
		Status Kamera & Berhasil \\
		Data Telemetri & Berhasil \\
		Notifikasi & Berhasil \\
		Data Gambar & Berhasil \\
		\hline
	\end{tabular}
\end{table}

Website berhasil memperbarui informasi tanpa memerlukan proses refresh halaman sehingga meningkatkan pengalaman pengguna dalam melakukan monitoring.

\section{Pengujian Motion Detection}

Pengujian dilakukan untuk mengetahui kemampuan sistem dalam mendeteksi aktivitas pada area pengawasan.

\subsection{Skenario Pengujian}

\begin{table}[H]
	\centering
	\caption{Skenario Pengujian Motion Detection}
	\begin{tabular}{|c|l|}
		\hline
		No & Kondisi \\
		\hline
		1 & Tidak terdapat pergerakan \\
		2 & Pergerakan kecil \\
		3 & Pergerakan sedang \\
		4 & Pergerakan besar \\
		\hline
	\end{tabular}
\end{table}

\subsection{Hasil Pengujian}

\begin{table}[H]
	\centering
	\caption{Hasil Pengujian Motion Detection}
	\begin{tabular}{|c|c|}
		\hline
		Kondisi & Hasil \\
		\hline
		Tidak Ada Gerakan & Tidak Terdeteksi \\
		Gerakan Kecil & Terdeteksi \\
		Gerakan Sedang & Terdeteksi \\
		Gerakan Besar & Terdeteksi \\
		\hline
	\end{tabular}
\end{table}

Hasil pengujian menunjukkan bahwa sistem mampu mendeteksi perubahan kondisi pada area pengawasan dan menghasilkan event yang digunakan untuk penyimpanan data serta pengiriman notifikasi.

\section{Pengujian Notifikasi Mobile}

Pengujian dilakukan untuk memastikan pengguna menerima notifikasi ketika terjadi aktivitas yang terdeteksi.

\subsection{Hasil Pengujian}

\begin{table}[H]
	\centering
	\caption{Hasil Pengujian Notifikasi Mobile}
	\begin{tabular}{|c|c|}
		\hline
		Parameter & Hasil \\
		\hline
		Pengiriman Notifikasi & Berhasil \\
		Penerimaan Notifikasi & Berhasil \\
		Navigasi ke Detail Kamera & Berhasil \\
		\hline
	\end{tabular}
\end{table}

Notifikasi berhasil diterima oleh aplikasi mobile dalam waktu yang relatif singkat setelah aktivitas terdeteksi.

\section{Pengujian Monitoring Telemetri}

Pengujian dilakukan untuk memastikan data kondisi perangkat dapat diterima dan ditampilkan pada dashboard monitoring.

\subsection{Parameter Telemetri}

\begin{enumerate}
	\item RSSI
	\item Uptime
	\item Free Heap
	\item Reset Reason
	\item Capture Success Rate
	\item Publish Success Rate
\end{enumerate}

\subsection{Hasil Pengujian}

\begin{table}[H]
	\centering
	\caption{Hasil Pengujian Telemetri}
	\begin{tabular}{|c|c|}
		\hline
		Parameter & Hasil \\
		\hline
		RSSI & Berhasil \\
		Uptime & Berhasil \\
		Free Heap & Berhasil \\
		Reset Reason & Berhasil \\
		Capture Success Rate & Berhasil \\
		Publish Success Rate & Berhasil \\
		\hline
	\end{tabular}
\end{table}

Data telemetri berhasil ditampilkan secara real-time sehingga memudahkan proses monitoring dan diagnosis perangkat.

\section{Pengujian Cloudflared Tunnel}

Pengujian dilakukan untuk memastikan sistem dapat diakses dari jaringan internet melalui domain yang telah dikonfigurasi.

\subsection{Hasil Pengujian}

\begin{table}[H]
	\centering
	\caption{Hasil Pengujian Cloudflared Tunnel}
	\begin{tabular}{|c|c|}
		\hline
		Parameter & Hasil \\
		\hline
		Akses Domain & Berhasil \\
		HTTPS & Berhasil \\
		Akses Dashboard & Berhasil \\
		API Backend & Berhasil \\
		\hline
	\end{tabular}
\end{table}

Cloudflared Tunnel berhasil menyediakan akses publik tanpa memerlukan pembukaan port secara langsung pada jaringan lokal.

\section{Pengujian Performa Server}

Pengujian dilakukan untuk mengetahui kemampuan server Intel NUC dalam menjalankan seluruh layanan secara bersamaan.

Parameter yang diamati meliputi:

\begin{enumerate}
	\item Penggunaan CPU
	\item Penggunaan RAM
	\item Penggunaan Storage
	\item Trafik Jaringan
\end{enumerate}

Hasil pengujian menunjukkan bahwa Intel NUC mampu menjalankan layanan Docker, Laravel, EMQX, MySQL, MinIO, WebSocket, dan Cloudflared secara bersamaan tanpa mengalami gangguan yang signifikan.

\section{Analisis Hasil}

Berdasarkan seluruh pengujian yang telah dilakukan, sistem berhasil memenuhi kebutuhan yang telah ditentukan pada tahap perancangan. Implementasi MQTT memungkinkan komunikasi data yang efisien antara perangkat dan server. Penggunaan WebSocket mampu menyediakan monitoring real-time dengan latensi yang lebih rendah dibandingkan mekanisme polling. Motion detection berhasil digunakan sebagai mekanisme deteksi aktivitas tanpa memerlukan komputasi yang tinggi seperti pada metode deteksi objek berbasis kecerdasan buatan.

Monitoring telemetri memberikan informasi kondisi perangkat secara berkala sehingga mempermudah proses identifikasi gangguan. Selain itu, integrasi notifikasi mobile memungkinkan pengguna memperoleh informasi aktivitas secara cepat. Penggunaan Intel NUC sebagai server utama juga mampu mendukung seluruh layanan yang dibutuhkan oleh sistem dengan performa yang baik dan stabil.